\newpage
\addcontentsline{toc}{section}{结束语}
\section*{结束语}

随着数据库在现代企业级应用中的核心地位日益凸显，确保数据库系统在各类变更和升级过程中的稳定性成为了运维领域的重要命题。本文着眼于数据库变更前评估的痛点，深入钻研了 PostgreSQL 的内核审计机制、事务时序模型与高并发调度技术，从零到一设计并实现了一套基于审计日志的高保真流量回放系统。

经过严谨的设计与开发，本文所实现的系统克服了海量乱序日志流解析、复杂事务上下文重构、以及微秒级极速时序对齐等一系列极具挑战性的工程难题。通过提出流式事务重构算法与基于滑动时间窗口的时序调整模型，成功突破了传统回放工具在请求保真度与回放性能间的瓶颈。综合评测不仅印证了系统在极限规模压力下依然能够维持优异的吞吐稳定性，更在时延分布、吞吐动态趋势等多个关键评估指标上全面展现了高达 98\% 以上的保真效果。

纵观整个研究历程，本文不仅在工程层面为 PostgreSQL 数据库生态输出了一套切实可行的流量回放与测试平台，也对“回放保真度”的量化评估方法进行了有益的学术探索。虽然现阶段的工作已能满足多数核心场景下的变更验证需求，但未来在多样化异构数据库支持与智能根因诊断领域，仍有广阔的发展空间。希望本文的研究能为业界同仁处理类似数据库流量评估与混沌工程问题提供有价值的参考。